\section{Функциональная декомпозиция, ссылки и константность}

\subsection{Как решать большую задачу так, чтобы решение было корректным, проверяемым, воспроизводимым, переиспользуемым и читаемым?}
\begin{definition}[Декомпозиция]
\textit{Декомпозиция} — это разбиение большой задачи на более маленькие, слабо связанные подзадачи, каждая из которых проще для понимания, реализации и проверки.
\end{definition}

Для разработки программ это означает:
\begin{itemize}
    \item разбивать решение на функции с ясной ответственностью;
    \item выделять интерфейс и реализацию;
    \item избегать дублирования логики;
    \item делать код тестируемым и воспроизводимым;
    \item использовать понятные имена и минимальные зависимости.
\end{itemize}

Хорошее решение обычно строится по схеме:
\begin{enumerate}
    \item формулируем входы, выходы и инварианты;
    \item выделяем подзадачи;
    \item реализуем каждую подзадачу отдельно;
    \item проверяем каждую часть независимо;
    \item затем связываем части в целое.
\end{enumerate}

\begin{remark}
Именно функциональная декомпозиция делает программу не просто «написанной», а инженерно организованной: её можно читать, тестировать, менять и переиспользовать.
\end{remark}

\subsection{Что такое функция с точки зрения языка C++? Как определить функцию или процедуру? Какие основные компоненты функции?}
\begin{definition}[Функция]
Функция в C++ — это именованный блок инструкций, который может получать параметры, выполнять вычисления и, возможно, возвращать результат.
\end{definition}

\begin{definition}[Процедура]
\textit{Процедура} — это функция, не возвращающая значение; в C++ это функция с типом возврата \texttt{void}.
\end{definition}

Функция обычно состоит из следующих частей:
\begin{itemize}
    \item \textbf{заголовок} — тип возврата, имя, список параметров;
    \item \textbf{тело} — блок инструкций;
    \item \textbf{инструкция \texttt{return}} — завершение функции и, если нужно, возврат значения.
\end{itemize}

\begin{example}
\begin{verbatim}
int sum(int a, int b) {
    return a + b;
}

void printHello() {
    std::cout << "Hello\n";
}
\end{verbatim}
\end{example}

\begin{remark}
Функция — это базовый строительный блок программы: через функции реализуется функциональная декомпозиция и уменьшается связанность кода.
\end{remark}

\subsection{Что такое неполный тип \texttt{void}? Можно ли породить объект такого типа? Что такое \texttt{void*}?}
\begin{definition}[void]
\texttt{void} — это неполный тип, обозначающий отсутствие значения или отсутствие конкретного типа.
\end{definition}

Объект типа \texttt{void} породить нельзя, потому что у него нет определённого размера и представления в памяти.

\texttt{void*} — это указатель на некий объект неопределённого типа. Такой указатель может хранить адрес, но сам по себе не позволяет разыменовывать его без приведения к конкретному типу.

\begin{itemize}
    \item \texttt{void} используется как тип возврата процедур;
    \item \texttt{void*} используется как универсальный указатель низкого уровня;
    \item в C++ \texttt{void*} применяется реже, чем в C, потому что C++ предпочитает типобезопасные абстракции.
\end{itemize}

\subsection{Свободные функции и функции-члены. Методы классов и статические функции-члены}
\begin{definition}[Свободная функция]
\textit{Free function} или \textit{non-member function} — функция, не принадлежащая классу.
\end{definition}

\begin{definition}[Функция-член]
\textit{Member function} — функция, объявленная внутри класса и работающая с его объектом.
\end{definition}

\begin{example}
\begin{verbatim}
int add(int a, int b) { return a + b; }   // free function

struct Foo {
    void bar() const;                     // member function
    static void baz();                    // static member function
};
\end{verbatim}
\end{example}

Методом называют функцию-член, обычно нестатическую.  
Статическая функция-член принадлежит классу как namespace-like сущность, а не конкретному объекту.

\begin{remark}
Свободная функция хороша как композиционная схема: она получает данные, обрабатывает их и возвращает результат без скрытого состояния. Для свободных функций обычно нежелательны побочные эффекты.
\end{remark}

\subsection{Частые уточнения: \texttt{std::getline}, статическая функция-член и методы}
\texttt{std::getline} — это \textbf{свободная функция} из стандартной библиотеки, а не метод \texttt{std::istream}.

Статическая функция-член (например, \texttt{Foo::baz()}) остаётся \textbf{функцией-членом}, но вызывается через имя класса:
\begin{verbatim}
Foo::baz();
\end{verbatim}

Метод — это нестатическая функция-член. Пример метода у \texttt{std::string}:
\begin{verbatim}
std::string s = "abc";
auto n = s.size(); // метод size()
\end{verbatim}

\subsection{Что такое side-effect и почему он нежелателен для свободных функций? Как \texttt{static} вносит состояние?}
\begin{definition}[Side effect]
\textit{Side effect} — побочный эффект, то есть изменение внешнего состояния: запись в поток, изменение глобальной переменной, модификация аргумента по ссылке, обращение к файлу и т.д.
\end{definition}

Побочный эффект нежелателен для свободной функции, потому что он:
\begin{itemize}
    \item ухудшает предсказуемость;
    \item усложняет тестирование;
    \item делает функцию менее переиспользуемой;
    \item затрудняет отладку.
\end{itemize}

\texttt{static} в теле функции позволяет сохранить состояние между вызовами:
\begin{verbatim}
int counter() {
    static int x = 0;
    return ++x;
}
\end{verbatim}

Здесь функция уже не полностью «чиста»: она зависит от предыдущих вызовов.

\subsection{Как устроено определение функции? Что такое prototype, caller и callee?}
\begin{definition}[Определение функции]
Определение функции — это описание её сигнатуры и тела.
\end{definition}

Обычно в определении есть:
\begin{itemize}
    \item имя функции;
    \item тип возврата;
    \item список формальных параметров;
    \item тело;
    \item при необходимости — \texttt{return}.
\end{itemize}

\begin{definition}[Prototype]
\textit{Prototype} — это объявление функции без тела, то есть её интерфейсная часть.
\end{definition}

\begin{definition}[Caller и callee]
\textit{Caller} — вызывающая функция; \textit{callee} — вызываемая функция.
\end{definition}

\begin{example}
\begin{verbatim}
int f(int);      // prototype
int g() {
    return f(10); // g is caller, f is callee
}
\end{verbatim}
\end{example}

\subsection{Список формальных параметров, заголовок и прототип}
Список формальных параметров описывает вход функции. В прототипе имена параметров необязательны:
\begin{verbatim}
int sum(int, int);      // корректный прототип
int sum(int a, int b);  // тоже корректный прототип
\end{verbatim}

Отличие:
\begin{itemize}
    \item \textbf{заголовок функции} — часть определения, за которой следует тело;
    \item \textbf{прототип} — объявление без тела, интерфейсная часть для других модулей трансляции.
\end{itemize}

\subsection{Функция: интерфейсная часть и реализация. Как это связано со сборкой программы?}
Интерфейс функции — это то, что видит пользователь функции:
\begin{itemize}
    \item имя;
    \item тип возврата;
    \item параметры;
    \item иногда спецификаторы, например \texttt{const} у member function.
\end{itemize}

Реализация — это тело функции.

Обычно распределение такое:
\begin{itemize}
    \item интерфейс — в заголовочном файле или модуле интерфейса;
    \item реализация — в \texttt{.cpp} или модуле реализации.
\end{itemize}

\begin{remark}
Компилятору часто достаточно интерфейса на этапе проверки вызова, но линкеру нужно полное определение функции, чтобы связать вызов с конкретным телом.
\end{remark}

\subsection{Точки выхода из функции: \texttt{return} в практике}
Функция может иметь несколько точек выхода через \texttt{return}. Это допустимо, если улучшает читаемость (например, ранний выход для проверки ошибок).

\begin{example}
\begin{verbatim}
int safeDiv(int a, int b) {
    if (b == 0) {
        return 0;
    }
    return a / b;
}
\end{verbatim}
\end{example}

Для \texttt{void}-функции инструкция \texttt{return;} также допустима и завершает выполнение процедуры.

\subsection{Признаки хорошей функции и принципы функциональной декомпозиции}
Хорошая функция обычно:
\begin{itemize}
    \item делает одну вещь;
    \item имеет ясное имя;
    \item имеет небольшой размер;
    \item имеет минимальное число параметров;
    \item по возможности не имеет побочных эффектов;
    \item легко тестируется;
    \item слабо связана с остальным кодом.
\end{itemize}

Принципы функциональной декомпозиции:
\begin{enumerate}
    \item одна функция — одна ответственность;
    \item сначала абстракция, потом детали;
    \item минимизировать связи между функциями;
    \item отделять вычисление от ввода-вывода;
    \item избегать дублирования логики.
\end{enumerate}

\begin{remark}
Симптомы плохой функции: слишком длинная, делает несколько разных задач, имеет много параметров, содержит много ветвлений, трудно тестируется и сложно читается.
\end{remark}

\subsection{Реентерабельность, параметризуемость, связность и тестируемость}
\begin{definition}[Реентерабельная функция]
Функция реентерабельна, если её можно безопасно вызывать повторно и параллельно без скрытого общего изменяемого состояния.
\end{definition}

\begin{definition}[Параметризуемость]
Параметризуемость — степень, в которой поведение функции задаётся входными параметрами, а не глобальным контекстом.
\end{definition}

Непараметризованный, но хороший пример: функция без входов, которая возвращает фиксированную доменную константу.
\begin{verbatim}
int maxScore() {
    return 100;
}
\end{verbatim}

Высокая связность означает, что шаги внутри функции решают одну задачу. Низкая связанность между функциями означает минимум лишних зависимостей.

Тестируемость растёт, когда функция:
\begin{itemize}
    \item имеет явные входы и выходы;
    \item не зависит от глобального состояния;
    \item не смешивает вычисления с вводом-выводом.
\end{itemize}

\subsection{Принципы S.O.L.I.D.}
\begin{definition}[S.O.L.I.D.]
\textit{S.O.L.I.D.} — набор принципов объектно-ориентированного и модульного проектирования.
\end{definition}

\begin{itemize}
    \item \textbf{S — Single Responsibility Principle}: у сущности должна быть одна причина для изменения.
    \item \textbf{O — Open/Closed Principle}: сущность должна быть открыта для расширения и закрыта для изменения.
    \item \textbf{L — Liskov Substitution Principle}: объекты наследника должны быть взаимозаменяемы с объектами базового типа.
    \item \textbf{I — Interface Segregation Principle}: лучше несколько узких интерфейсов, чем один широкий.
    \item \textbf{D — Dependency Inversion Principle}: зависимости должны строиться на абстракциях, а не на конкретных реализациях.
\end{itemize}

\subsection{Формальные и фактические параметры. Сигнатура и перегрузка функций}
\begin{definition}[Формальные и фактические параметры]
\textit{Формальные параметры} — переменные в объявлении функции.  
\textit{Фактические параметры} — значения, передаваемые при вызове.
\end{definition}

\begin{example}
\begin{verbatim}
int sum(int a, int b);   // a, b — formal parameters
int x = sum(2, 3);       // 2, 3 — actual arguments
\end{verbatim}
\end{example}

\begin{definition}[Сигнатура функции]
Сигнатура — это характеристика функции, по которой различаются перегрузки: имя и типы параметров; для member functions часто также учитываются cv-квалификаторы.
\end{definition}

\begin{definition}[Перегрузка]
\textit{Overload} — несколько функций с одним именем, но различающихся сигнатурой.
\end{definition}

\begin{remark}
Перегрузка не может отличаться только возвращаемым типом: для выбора нужной функции важны имя и типы параметров, а не результат. Также значения по умолчанию не создают отдельную перегрузку.
\end{remark}

\begin{remark}
Сигнатура нужна отдельно от полного заголовка, чтобы различать именно «форму вызова», а не детали реализации.
\end{remark}

\subsection{Модели взаимодействия функций}
Взаимодействие функций удобно изображать:
\begin{itemize}
    \item графом вызовов функций;
    \item UML sequence diagram;
    \item схемой потоков данных.
\end{itemize}

\begin{remark}
Граф вызовов показывает структуру зависимости, а sequence diagram — порядок сообщений и переходов управления во времени.
\end{remark}

\subsection{CV-квалификаторы. \texttt{const} и \texttt{volatile}}
\begin{definition}[CV-квалификаторы]
\textit{CV-qualified} тип — это тип с квалификаторами \texttt{const} и/или \texttt{volatile}.
\end{definition}

\texttt{const} означает, что объект нельзя изменять через данное имя.  
\texttt{volatile} означает, что значение может измениться вне контроля программы, поэтому компилятор не должен агрессивно оптимизировать обращения.

\begin{example}
\begin{verbatim}
const int a = 5;
int const b = 5;
\end{verbatim}
\end{example}

\texttt{const int} и \texttt{int const} эквивалентны: различается только синтаксическая запись, а не смысл.

\begin{remark}
\texttt{const} используется часто, \texttt{volatile} — редко, в основном для низкоуровневого взаимодействия с оборудованием или сигналами извне.
\end{remark}

\subsection{Константный объект: что сохраняется и что запрещается}
Объект, объявленный через \texttt{const}, остаётся полноценным объектом языка: у него есть имя (если оно задано), адрес, тип и значение.

\texttt{const} влияет именно на операционную семантику: через данное имя нельзя выполнять модифицирующие операции над объектом.

\begin{example}
\begin{verbatim}
const int limit = 100;
int x = limit;  // читать можно
// limit = 200; // нельзя: ошибка компиляции
\end{verbatim}
\end{example}

Это не делает объект «нематериальным»: память у него есть, и адрес можно получить (при соблюдении правил типа и времени жизни).

\subsection{Ссылка как ещё одно имя для объекта}
\begin{definition}[Ссылка]
\textit{Reference} — это альтернативное имя для уже существующего объекта.
\end{definition}

Объект может иметь:
\begin{itemize}
    \item 0 имен, если доступен только через адрес;
    \item 1 имя, если существует обычная переменная;
    \item много имён, если есть ссылки, псевдонимы, параметры по ссылке и т.д.
\end{itemize}

\begin{remark}
Мнемоника такая: «объекту нужно ещё одно имя, если старое имя недоступно, неудобно или его нужно передать в другой контекст без копирования».
\end{remark}

\subsection{Создание ссылки, binding, перепривязка, отложенное привязывание}
Ссылка создаётся через \texttt{\&} в объявлении:
\begin{verbatim}
int x = 10;
int& r = x;
\end{verbatim}

Это называется \textit{binding}: ссылка привязывается к объекту.

\begin{remark}
Обычную ссылку перепривязать нельзя. После привязки она остаётся связанной с тем же объектом до конца жизни.
\end{remark}

Отложенное привязывание обычной ссылки невозможно.  
Если нужно «позже связать» сущность с объектом, используют указатель, \texttt{std::optional<std::reference\_wrapper<T>>} или другой уровень абстракции.

\subsection{Ссылки в range-based for}
Ссылка в range-based for нужна, когда нужно модифицировать элементы диапазона без копирования:
\begin{verbatim}
for (auto& x : v) {
    x *= 2;
}
\end{verbatim}

\begin{remark}
Без \texttt{\&} создавалась бы копия элемента, и изменения не повлияли бы на исходную коллекцию.
\end{remark}

\subsection{Ссылки как параметры функций}
Передача по ссылке позволяет callee изменить объект caller:
\begin{verbatim}
void inc(int& x) {
    ++x;
}
\end{verbatim}

\begin{definition}[Ссылка как двунаправленный портал]
Неконстантная ссылка позволяет функции «видеть» и изменять объект caller напрямую.
\end{definition}

Использование:
\begin{itemize}
    \item для модификации аргумента;
    \item для возврата нескольких значений через параметры;
    \item для избежания копирования больших объектов.
\end{itemize}

Пример возврата нескольких результатов без агрегатов:
\begin{verbatim}
void minMax(int a, int b, int& mn, int& mx) {
    if (a < b) {
        mn = a;
        mx = b;
    } else {
        mn = b;
        mx = a;
    }
}
\end{verbatim}

\subsection{Передача по значению и по ссылке. Правило по умолчанию}
\begin{remark}
В C++ по умолчанию всё передаётся по значению: копируется.
\end{remark}

Передача по значению:
\begin{itemize}
    \item хороша для малых объектов;
    \item защищает исходный объект от изменения;
    \item создаёт копию, что может быть дорого.
\end{itemize}

Передача по ссылке:
\begin{itemize}
    \item удобна для больших объектов;
    \item избегает копирования;
    \item позволяет изменять объект;
    \item требует аккуратности с временем жизни объекта.
\end{itemize}

\subsection{Константная ссылка}
\begin{definition}[Константная ссылка]
\textit{Const reference} — это ссылка только для чтения: через неё нельзя изменять объект.
\end{definition}

Применение:
\begin{itemize}
    \item передача больших объектов без копирования;
    \item защита от случайного изменения;
    \item чтение элементов в range-based for:
\end{itemize}

\begin{example}
\begin{verbatim}
for (const auto& x : v) {
    std::cout << x << '\n';
}
\end{verbatim}
\end{example}

\begin{remark}
Константная ссылка особенно полезна для «больших» объектов. Для «малых» копирование часто дешевле, чем усложнение логики ссылок.
\end{remark}

\subsection{Что такое «большие» и «малые» объекты? Правило трёх указателей}
\begin{definition}[Правило трёх указателей]
Объект считается «большим», если его копирование потенциально дорого и он по размеру сравним с тремя указателями или более, либо если он владеет динамическим ресурсом.
\end{definition}

Неформально:
\begin{itemize}
    \item \textbf{малые объекты}: \texttt{int}, \texttt{double}, \texttt{char}, небольшие \texttt{enum};
    \item \textbf{большие объекты}: \texttt{std::string}, контейнеры STL, крупные структуры, объекты с динамической памятью.
\end{itemize}

\subsection{Идиоматические правила передачи объектов в функции}
Обычно используют такие правила:
\begin{enumerate}
    \item \textbf{большие объекты для чтения} — через \textbf{константную ссылку};
    \item \textbf{малые объекты для чтения} — \textbf{по значению};
    \item \textbf{объекты для чтения и записи} — через \textbf{неконстантную ссылку}.
\end{enumerate}

\begin{example}
\begin{verbatim}
void print(const std::string& s);  // read-only large object
int square(int x);                 // small object by value
void normalize(std::vector<double>& v); // read/write by reference
\end{verbatim}
\end{example}

\subsection{Краткие ответы на контрольные вопросы}
\begin{itemize}
    \item \textbf{Что такое функция?} Именованный блок инструкций с входом и, возможно, выходом.
    \item \textbf{Что такое \texttt{void}?} Неполный тип, обозначающий отсутствие значения.
    \item \textbf{Что такое свободная функция?} Не член класса.
    \item \textbf{Что такое статическая функция-член?} Функция, принадлежащая классу, но не конкретному объекту.
    \item \textbf{Что такое \texttt{const}?} Запрет на изменение через данное имя.
    \item \textbf{Что такое ссылка?} Ещё одно имя для объекта.
    \item \textbf{Можно ли перепривязать ссылку?} Нет.
    \item \textbf{Когда передавать по ссылке?} Когда нужен доступ к исходному объекту или надо избежать копирования.
    \item \textbf{Когда передавать по значению?} Когда объект маленький или копия дешева и нужна изоляция.
    \item \textbf{Когда нужна константная ссылка?} Когда объект большой, но его надо только читать.
\end{itemize}
